% Abstract (included inside the abstract environment of main_phase2.tex, so no
% \section). IEEE rule: no citations, footnotes, math or special symbols, so
% "percent" is written out. Numbers: results/phase2/metrics/comparison.csv,
% ablation.csv, seeds.csv and tuning/blend.json; Phase 1: results/metrics/.
GPS fails inside buildings, but the Wi-Fi signals a phone hears show where it
is. On the public UJIIndoorLoc data we predict a scan's building, floor and
position with one multi-task neural network: shared hidden layers feed a
building head, a floor head and a position head. Its inputs are the scaled
signal strengths and heard indicators of the access points, plus our Phase 1
features as a first hybrid. As in Phase 1, settings are tuned by cross-validation
that holds out whole users with their phones, and the test set, recorded months
later mostly with new phones, is used only at the end. An ablation study shows
that dropout with weight decay and the heard indicators help, while the Phase 1
features do not, and that three single-task networks beat the shared one. In
cross-validation the network beat our Phase 1 Random Forest (9.21 against 11.09
m mean 2D error), but on the test set it reached 8.32 m and 90.8 percent of
floors, against 7.42 m and 92.2 percent for the forest; over five seeds it
scored 8.21 m. A second hybrid, a blend of the forest and the network with
one weight per task chosen by cross-validation, beat both in cross-validation
(9.14 m) and improved the network on the test set (8.10 m and 92.2 percent of
floors), but not the forest's position error.
